% ============================================================================
%  CHAPITRE 3 — SECTIONS COMPLÉTÉES AVEC MESURES RÉELLES
%
%  Toutes les valeurs de ce fichier ont été relevées sur la machine de
%  développement le 17 août 2026, soit par interrogation du système, soit par
%  exécution effective des suites de tests, soit par comptage sur le dépôt.
%  Les emplacements restés en [à compléter] exigent la pile complète démarrée.
% ============================================================================


% ============================================================================
%  REMPLACE : tableau tab:env-materiel
% ============================================================================

\begin{table}[htbp]
\centering
\caption{Environnement d'implémentation et d'évaluation}
\label{tab:env-materiel}
\small
\begin{tabularx}{\textwidth}{|p{4.4cm}|X|}
\hline
\textbf{Composant} & \textbf{Spécification} \\ \hline

Système d'exploitation & Microsoft Windows 11 Éducation, version 10.0.26200 \\ \hline
Processeur & Intel Core i5-12500H, douzième génération : 12 cœurs physiques,
16 fils d'exécution, fréquence de base 2,5 GHz \\ \hline
Mémoire vive & 15,6 Gio \\ \hline
Stockage & Micron 2400 NVMe, 477 Gio \\ \hline
Topologie de déploiement & Ensemble des services sur un hôte unique \\ \hline
Java & 17.0.12 LTS pour l'interface métier et Spark, 21 embarqué pour Hop \\ \hline
Spring Boot & 3.4.5 \\ \hline
Apache Kafka & 8.2.2 en mode KRaft \\ \hline
Apache Hop & 2.18.0 \\ \hline
Apache Spark & 3.5.7 \\ \hline
PostgreSQL & 16 \\ \hline
MongoDB & 8.2 \\ \hline
RustFS & 1.0.0-beta.12 \\ \hline
HashiCorp Vault & 1.21.7 \\ \hline
Keycloak & haute disponibilité, deux nœuds \\ \hline
ClamAV & 1.4 \\ \hline
Nginx & 1.29-alpine \\ \hline
Docker & 27.1.1 avec Compose v2 \\ \hline
Node.js & 20.16.0 pour le portail et le médiateur \\ \hline
Maven & 3.9.9 \\ \hline
Jeu d'essai principal & \emph{[à compléter : nature, volume, période couverte]} \\ \hline
Sources sollicitées & \emph{[à compléter : types et nombre]} \\ \hline

\end{tabularx}
\end{table}

\noindent
La mémoire disponible mérite une remarque, car elle conditionne la lecture des
mesures à venir. Les 15,6 Gio de l'hôte sont partagés entre l'ensemble des
services, dont trois sont notoirement consommateurs : le moteur distribué, le
serveur d'identité et l'antivirus. Cette contrainte a d'ailleurs orienté une
décision d'implémentation, l'interface métier ayant reçu un plafond mémoire
explicite après qu'une extraction volumineuse eut épuisé les ressources de la
machine.


% ============================================================================
%  AJOUT : nouvelle sous-section à insérer après « Organisation du dépôt »
% ============================================================================

\subsection{Volumétrie du code produit}

Le tableau \ref{tab:volumetrie-code} donne la taille des modules développés. Cette
volumétrie n'est pas un indicateur de qualité, mais elle situe l'ampleur de la
réalisation et la répartition de l'effort entre les couches.

\begin{table}[htbp]
\centering
\caption{Volumétrie du code produit par module}
\label{tab:volumetrie-code}
\small
\begin{tabularx}{\textwidth}{|p{3.6cm}|p{2.2cm}|p{2cm}|X|}
\hline
\textbf{Module} & \textbf{Lignes} & \textbf{Fichiers} & \textbf{Langage} \\ \hline

\texttt{api-core} & 16 970 & 156 & Java \\ \hline
\texttt{frontend} & 14 416 & --- & TypeScript \\ \hline
\texttt{source-\allowbreak gateway} & 6 250 & 36 & Java \\ \hline
\texttt{pipeline-\allowbreak consumer} & 3 160 & 13 & Java \\ \hline
\texttt{iol-mediator} & 3 154 & --- & JavaScript \\ \hline
Moteurs d'exécution & 1 156 & 3 & Python \\ \hline
\textbf{Total} & \textbf{45 106} & & \\ \hline

\end{tabularx}
\end{table}

\noindent
La part de l'interface métier, qui concentre plus du tiers du code, s'explique par
le nombre de responsabilités qu'elle porte : exposition des interfaces,
authentification, gestion des métadonnées, assemblage du contrat d'exécution et
validation du SQL. La passerelle de sources, extraite en cours de réalisation,
représente à elle seule un quart du code Java, ce qui donne la mesure du plan de
données séparé du plan de contrôle.


% ============================================================================
%  REMPLACE : sous-section « Campagne de tests automatisés » en entier
% ============================================================================

\subsection{Campagne de tests automatisés}

Les services du système comprennent deux cent cinquante-trois tests automatisés.
La totalité a été exécutée avec succès, sans échec ni test ignoré, sur
l'environnement décrit au tableau \ref{tab:env-materiel}.

\begin{table}[htbp]
\centering
\caption{Répartition des tests automatisés par module}
\label{tab:tests-modules}
\small
\begin{tabularx}{\textwidth}{|p{3.2cm}|p{1.4cm}|p{2cm}|X|}
\hline
\textbf{Module} & \textbf{Nombre} & \textbf{Résultat} & \textbf{Domaines couverts} \\ \hline

\texttt{api-core} & 134 & 134 réussis & Soumission asynchrone, sécurité, audit, contrat de
métadonnées, validation du SQL, garde de confidentialité, idempotence, sanitisation \\ \hline
\texttt{source-\allowbreak gateway} & 38 & 38 réussis & Acquittement, réservation, purge des identifiants,
compatibilité cryptographique, parité du modèle de lecture \\ \hline
\texttt{pipeline-\allowbreak consumer} & 30 & 30 réussis & Orchestration des moteurs, verrouillage
d'exécution, gestion des filigranes \\ \hline
\texttt{iol-mediator} & 51 & 51 réussis & Interopérabilité, idempotence, garde contre les
requêtes falsifiées \\ \hline
\textbf{Total} & \textbf{253} & \textbf{253 réussis} & \\ \hline

\end{tabularx}
\end{table}

Certains tests ne portent pas sur des fonctions isolées mais sur les propriétés
critiques de la conception. Le tableau \ref{tab:tests-critiques} recense ceux dont
l'échec signalerait une atteinte à une garantie architecturale, avec le nombre de
cas que chacun couvre.

\begin{table}[htbp]
\centering
\caption{Tests portant sur les propriétés critiques de la conception}
\label{tab:tests-critiques}
\small
\begin{tabularx}{\textwidth}{|p{4.4cm}|p{1.4cm}|X|}
\hline
\textbf{Suite} & \textbf{Cas} & \textbf{Propriété garantie} \\ \hline

Validation de sûreté du SQL & 18 & Refus des suppressions, des modifications de
structure et des attributions de droits \\ \hline
Parité de sanitisation & 12 & Identité stricte entre la normalisation Java et son
équivalent Python \\ \hline
Transport et purge & 9 & Aucun identifiant de source ne subsiste dans la commande
publiée \\ \hline
Résolution de la destination & 10 & Injection correcte de la connexion cible dans
chaque source \\ \hline
Filigranes incrémentaux & 6 & Persistance de la borne de reprise par source \\ \hline
Réservation d'exécution & 7 & Absence de double extraction après rééquilibrage \\ \hline
Soumission asynchrone & 5 & Refus d'une exécution concurrente et de la saturation \\ \hline
Idempotence de réception & 4 & Un message reçu deux fois n'est traité qu'une fois \\ \hline
Garde de confidentialité & 2 & Aucune donnée métier transmise à un fournisseur
externe \\ \hline

\end{tabularx}
\end{table}

Deux vérifications complètent ce dispositif sans relever des tests unitaires. La
première est un ensemble de dix-sept invariants évalués au démarrage de
l'interface métier en profil de production : le service refuse de démarrer si le
chiffrement transitoire, l'authentification déléguée, le transport chiffré ou
l'analyse antivirale en mode fermé ne sont pas actifs. La seconde est une porte de
sécurité comptant quatre-vingt-seize contrôles sur la topologie résolue, exécutée
en intégration continue, qui vérifie notamment la réplication du bus et des bases,
l'absence de port publié hors du mandataire inverse et le caractère immuable des
images déployées.


% ============================================================================
%  REMPLACE : sous-section « Couverture des protocoles » (question 1)
% ============================================================================

\subsection{Couverture des protocoles}

Le moteur universel déclare quinze types de sources. Leur décompte appelle une
distinction, car ils ne relèvent pas tous du même mécanisme. Quatorze
correspondent à des protocoles d'accès que la plateforme sollicite elle-même :
huit systèmes relationnels, cinq formats de fichier et une interface HTTP. Le
quinzième, le dépôt entrant, n'est pas un protocole d'accès mais un canal par
lequel un partenaire pousse ses données vers la plateforme ; il est pris en charge
par la couche d'interopérabilité et non par le moteur de lecture.

\begin{table}[htbp]
\centering
\caption{Types de sources déclarés et état de leur validation}
\label{tab:protocoles}
\small
\begin{tabularx}{\textwidth}{|p{3.2cm}|p{1.6cm}|p{4cm}|X|}
\hline
\textbf{Famille} & \textbf{Déclarés} & \textbf{Détail} & \textbf{État de validation} \\ \hline

Bases relationnelles & 8 & PostgreSQL, MySQL, MariaDB, SQL Server, Oracle, SQLite,
Snowflake, Redshift & Scénario S1 à consigner \\ \hline
Formats de fichier & 5 & CSV, Excel, Parquet, Avro, ORC & Scénario S2 à consigner \\ \hline
Interface HTTP & 1 & API REST au format JSON & Scénario S3 à consigner \\ \hline
Dépôt entrant & 1 & Canal d'interopérabilité & Scénarios S6 et S7 \\ \hline
\textbf{Total} & \textbf{15} & & \\ \hline

\end{tabularx}
\end{table}

Les huit pilotes relationnels sont effectivement embarqués dans la passerelle, ce
qui a été vérifié dans sa déclaration de dépendances. La cible de plus de quinze
protocoles fixée au chapitre précédent est donc atteinte si l'on compte le canal
entrant, et manquée d'une unité si l'on s'en tient aux protocoles d'accès. Nous
retenons la seconde lecture, plus exigeante : elle seule mesure ce que l'hypothèse
H1 énonce, à savoir l'indépendance du pipeline vis-à-vis de la nature des sources
qu'il va lui-même interroger.

Cette couverture reste par ailleurs déclarative. Elle établit l'étendue visée par
le mécanisme, non le fait que chaque protocole ait été éprouvé de bout en bout.


% ============================================================================
%  REMPLACE : tableau tab:res-h1
% ============================================================================

\begin{table}[htbp]
\centering
\caption{Indicateurs relatifs à l'hypothèse H1}
\label{tab:res-h1}
\small
\begin{tabularx}{\textwidth}{|p{4.6cm}|p{2.6cm}|X|}
\hline
\textbf{Indicateur} & \textbf{Cible} & \textbf{Mesure} \\ \hline

Lignes de code à écrire pour intégrer une source nouvelle & Aucune & Aucune :
l'intégration se réduit à une configuration \\ \hline
Protocoles d'accès pris en charge & Supérieur à quinze & Quatorze protocoles
d'accès déclarés, quinze types de sources au total ; cible non atteinte selon la
lecture retenue \\ \hline
Pilotes relationnels embarqués & --- & Huit, vérifiés dans la déclaration de
dépendances de la passerelle \\ \hline
Redéploiement nécessaire & Aucun & Aucun : la prise en compte est dynamique \\ \hline
Durée de configuration d'un flux & Non fixée & \emph{[à compléter : mesure sur
scénario S1]} \\ \hline

\end{tabularx}
\end{table}


% ============================================================================
%  REMPLACE : sous-section « Bilan de l'hypothèse H1 »
% ============================================================================

\subsection{Bilan de l'hypothèse H1}

Le mécanisme attendu par H1 est implémenté : l'intégration d'une source repose sur
le contrat de métadonnées et ne requiert ni code spécifique ni redéploiement.
Quatorze protocoles d'accès sont déclarés, dont huit pilotes relationnels dont la
présence effective a été vérifiée. La validation demeure toutefois partielle sur
deux points. La cible de plus de quinze protocoles n'est pas atteinte au sens
strict, et les scénarios S1 à S3 doivent encore consigner les protocoles
effectivement éprouvés sur une source représentative. H1 ne sera considérée comme
pleinement validée qu'après cette campagne.


% ============================================================================
%  REMPLACE : sous-section « Ajout d'une norme par déclaration » (question 2)
% ============================================================================

\subsection{Ajout d'une norme par déclaration}

Les trois normes sectorielles actuellement configurées, FHIR R4, ISO 20022 et
Ed-Fi, sont prises en charge par des médiateurs dédiés qui encodent une partie de
leurs règles. Les cinquante et un tests du médiateur générique, tous réussis,
couvrent la validation contre le référentiel, l'idempotence de la réception et la
protection contre les requêtes falsifiées vers des adresses internes.

Ces éléments démontrent la capacité d'interopérabilité sectorielle de la
plateforme, mais ne suffisent pas à établir la propriété d'extension sans code du
médiateur générique. Cette propriété devra être éprouvée par l'ajout d'une norme
nouvelle au seul moyen du référentiel, sans modification ni redéploiement d'aucun
service.


% ============================================================================
%  REMPLACE : sous-section « Couverture de la journalisation » (question 4)
% ============================================================================

\subsection{Couverture de la journalisation}

Le journal d'audit enregistre neuf catégories d'action, définies par
l'énumération correspondante du domaine : création, lecture, modification,
suppression, exécution, activation, dépublication, export et import. Chaque entrée
consigne l'auteur, son rôle, le type et l'identifiant de la ressource concernée,
l'issue de l'opération, l'adresse réseau d'origine et la durée du traitement.

La couverture effective de ces catégories sur l'ensemble des points d'entrée
d'administration reste à consigner par un relevé sur la plateforme en
fonctionnement.


% ============================================================================
%  REMPLACE : tableau tab:res-h4
% ============================================================================

\begin{table}[htbp]
\centering
\caption{Indicateurs relatifs à l'hypothèse H4}
\label{tab:res-h4}
\small
\begin{tabularx}{\textwidth}{|p{4.6cm}|p{2.6cm}|X|}
\hline
\textbf{Indicateur} & \textbf{Cible} & \textbf{Mesure} \\ \hline

Catégories d'action journalisées & Toute action d'administration & Neuf catégories
définies ; couverture des points d'entrée \emph{[à compléter]} \\ \hline
Invariants de sécurité au démarrage & --- & Dix-sept, bloquant le démarrage en
production s'ils ne sont pas satisfaits \\ \hline
Contrôles sur la topologie déployée & --- & Quatre-vingt-seize, évalués en
intégration continue \\ \hline
Motifs de rejet du garde de confidentialité & --- & Neuf, appliqués avant tout
envoi à un fournisseur externe \\ \hline
Reconstitution de l'historique d'une donnée & De la réception à l'exposition &
\emph{[à compléter : scénario à conduire]} \\ \hline
Absence d'identifiant source hors de la passerelle & Vérifiée & Partielle :
l'interface métier déchiffre encore des identifiants pour la découverte de schéma,
l'atelier SQL, le test de connexion, l'estimation et le chemin entrant \\ \hline
Refus d'un message contenant une source directe & Systématique & Systématique :
assuré par l'assertion récursive de purge, couverte par neuf cas de test \\ \hline

\end{tabularx}
\end{table}


% ============================================================================
%  REMPLACE : sous-section « Cadrage des résultats présentés », second paragraphe
% ============================================================================

Le second niveau porte sur les \textbf{mesures de performance} : durées
d'exécution, volumes traités et latences. La campagne correspondante doit être
conduite sur le poste décrit au tableau \ref{tab:env-materiel}, où quinze services
partagent 15,6 Gio de mémoire et seize fils d'exécution. Les valeurs qui en
résulteront caractériseront cet environnement et ne préjugeront pas du
comportement sur une infrastructure dimensionnée. Cette réserve n'est pas de
principe : elle découle d'une observation faite en cours de réalisation, où la
concurrence entre le portail et une extraction volumineuse sur les mêmes
ressources a conduit à séparer le plan de données du plan de contrôle.


% ============================================================================
%  REMPLACE : tableau tab:synthese-hypotheses
% ============================================================================

\begin{table}[htbp]
\centering
\caption{Synthèse du statut des quatre hypothèses}
\label{tab:synthese-hypotheses}
\small
\begin{tabularx}{\textwidth}{|p{0.8cm}|p{3.6cm}|p{3cm}|X|}
\hline
& \textbf{Hypothèse} & \textbf{Statut} & \textbf{Réserve} \\ \hline

H1 & Généricité de l'intégration & Implémentée, validation partielle & Quatorze
protocoles d'accès déclarés pour une cible de quinze ; scénarios S1 à S3 à
consigner \\ \hline
H2 & Interopérabilité sans code & Implémentée, validation à conduire & Ajout d'une
norme par le médiateur générique et échange bidirectionnel à démontrer \\ \hline
H3 & Qualité et rejouabilité & Implémentée, validation à conduire & Rejeu mesuré à
effectuer ; entrepôt sur moteur mono-instance \\ \hline
H4 & Sécurité et traçabilité & Implémentée, validation partielle & Historique de
bout en bout à reconstituer et confinement des identifiants encore partiel \\ \hline

\end{tabularx}
\end{table}


% ============================================================================
%  REMPLACE : Conclusion du chapitre
% ============================================================================

\section*{Conclusion}
\addcontentsline{toc}{section}{Conclusion}

En définitive, la réalisation présentée dans ce chapitre établit la faisabilité
technique d'une plateforme d'intégration pilotée par métadonnées, combinant
raffinage des données, interopérabilité, sécurité et supervision. Les mécanismes
associés aux quatre hypothèses sont effectivement implémentés, et les deux cent
cinquante-trois tests automatisés, tous exécutés avec succès, apportent un premier
niveau de validation portant notamment sur les propriétés critiques de la
conception : sûreté du SQL, purge des identifiants, idempotence, verrouillage
d'exécution et parité de sanitisation.

Cette validation reste néanmoins d'ordre unitaire. Sa confirmation au niveau du
système exige l'exécution des scénarios de bout en bout et le relevé des mesures
de volume et de durée qui leur sont associées. Le chapitre distingue ainsi
délibérément deux natures de résultats : les mécanismes établis par le code et ses
tests, d'une part, et les propriétés qui ne se démontrent que par une exécution
complète sur données représentatives, d'autre part. Cette distinction, plutôt
qu'une présentation indifférenciée, nous paraît la condition d'une évaluation
honnête du travail accompli.
